% GENERATED FILE. Edit canonical lessons, then run npm run generate:books.
% canonical-source-hash: fnv1a64:bb5f4219b0ae63a9
% canonical-lessons: ES-C16-imperfecto, ES-C16-comer-imperfecto, ES-C16-vivir-imperfecto, ES-C16-ver, ES-C16-ser-imperfecto, ES-C16-ir-imperfecto, ES-C16-ver-imperfecto, ES-C16-practice

\chapter{The Imperfect --- the Past That Kept Going}
\label{ch:imperfect}
\hlchaptermodality{37}

\begin{chapteropening}
\textbf{By the end of this chapter:} \emph{I can use the Spanish imperfect for past habits, background, and actions in progress.}

Complete ES-C16-practice and retrieve the bounded singular imperfect rows for hablar, comer, vivir, ser, ir, and newly learned ver while contrasting open and completed past viewpoints.
\end{chapteropening}

\section[hablaba]{hablaba — let a past action stay open}

\label{lesson:ES-C16-imperfecto}



\begin{quote}
\textbf{Hablé español} presents speaking as a completed past event. A different past, the \textbf{imperfect}, can leave the action open: \textbf{Hablaba español} — “I was speaking Spanish” or “I used to speak Spanish.” The sentence does not say where the action ended.
\end{quote}

\begin{grammarlens}[title={Grammar Lens: three familiar persons}]
\noindent
\begin{tabularx}{\linewidth}{@{}>{\raggedright\arraybackslash}X>{\raggedright\arraybackslash}X@{}}
\toprule
\textbf{person} & \textbf{open or habitual past} \\
\midrule
yo & habl\textbf{aba} \\
tú & habl\textbf{abas} \\
él / ella / usted & habl\textbf{aba} \\
\bottomrule
\end{tabularx}

Say \textbf{hablaba · hablabas · hablaba}. The first and third forms are identical; the surrounding conversation or an explicit person tells you who is speaking. Plural persons wait.
\end{grammarlens}

\begin{cousinweb}
Spanish \textbf{-aba} continues the Latin imperfect family in \textbf{-ābam}. Historical sound change reshaped the complete Latin forms, but the \textbf{b} remains a useful audible bridge between \textbf{-ābam} and Spanish \textbf{-aba}.
\end{cousinweb}

\subsection*{Your turn}
Say these aloud:

\begin{itemize}
\raggedright
  \item “hablaba, hablabas, hablaba”
  \item “Hablé español. Hablaba español.”
  \item “-ābam and -aba share b”
\end{itemize}

\begin{tcolorbox}[breakable,colback=teal!4,colframe=teal!35!black,title={Before you move on}]
Give the three learned forms. (\emph{Hablaba, hablabas, hablaba.}) What can the imperfect do? (Leave a past action ongoing or habitual.) Which Latin family stands behind \textbf{-aba}? (\textbf{-ābam}.) Next: build the same gentle frame with \textbf{comer}.
\end{tcolorbox}
\section[comía]{comía — hear both vowels}

\label{lesson:ES-C16-comer-imperfecto}



\begin{quote}
\textbf{Comí} presents eating as completed. \textbf{Comía} leaves it open or habitual: “I was eating” or “I used to eat.” One accent separates the two forms in writing and in sound.
\end{quote}

\begin{grammarlens}[title={Grammar Lens: the singular -er row}]
\noindent
\begin{tabularx}{\linewidth}{@{}>{\raggedright\arraybackslash}X>{\raggedright\arraybackslash}X@{}}
\toprule
\textbf{person} & \textbf{open or habitual past} \\
\midrule
yo & com\textbf{ía} \\
tú & com\textbf{ías} \\
él / ella / usted & com\textbf{ía} \\
\bottomrule
\end{tabularx}

Say \textbf{comía · comías · comía}. As with \textbf{hablaba}, the first and third forms match. Plural persons wait.
\end{grammarlens}

\begin{grammarlens}[title={Grammar Lens: co-mí-a}]
The accent in \textbf{-ía} keeps \textbf{i} and \textbf{a} in separate syllables: \textbf{co-mí-a}, three beats. It appears in every form learned here. Compare final stress in \textbf{comí} with the middle stress and extra syllable in \textbf{comía}.
\end{grammarlens}

\begin{cousinweb}
The \textbf{-er} imperfect continues a Romance reshaping of Latin \textbf{-ēbam} material. Among the changes, an intervocalic \textbf{b} was lost; later Spanish spelling marks the stressed \textbf{í} that keeps \textbf{i-a} apart. This is a layered history, not a single instruction to delete one letter.
\end{cousinweb}

\subsection*{Your turn}
Say these aloud:

\begin{itemize}
\raggedright
  \item “comía, comías, comía”
  \item “Comí. Comía.”
  \item “co-mí-a”
\end{itemize}

\begin{tcolorbox}[breakable,colback=teal!4,colframe=teal!35!black,title={Before you move on}]
Give the three learned forms. (\emph{Comía, comías, comía.}) What does the accent do? (Keeps \textbf{i} and \textbf{a} apart.) Which Latin family lies behind the ending? (\textbf{-ēbam}, reshaped through Romance.) Next: test whether \textbf{vivir} can reuse this row.
\end{tcolorbox}
\section[vivía]{vivía — reuse the complete -er row}

\label{lesson:ES-C16-vivir-imperfecto}



\begin{quote}
Chapter 15 let \textbf{comer} and \textbf{vivir} share one completed-past row. They converge again here. Replace \textbf{com-} in \textbf{comía, comías, comía} with \textbf{viv-}.
\end{quote}

\begin{grammarlens}[title={Grammar Lens: one row reused}]
\noindent
\begin{tabularx}{\linewidth}{@{}>{\raggedright\arraybackslash}X>{\raggedright\arraybackslash}X>{\raggedright\arraybackslash}X@{}}
\toprule
\textbf{person} & \textbf{comer} & \textbf{vivir} \\
\midrule
yo & comía & vivía \\
tú & comías & vivías \\
él / ella / usted & comía & vivía \\
\bottomrule
\end{tabularx}

Say \textbf{vivía · vivías · vivía}. \textbf{-er} and \textbf{-ir} share \textbf{-ía, -ías, -ía}. The accent still keeps \textbf{i-a} apart. Plural persons wait.
\end{grammarlens}

\begin{grammarlens}[title={Grammar Lens: one known contrast}]
\begin{itemize}
\raggedright
  \item \textbf{Viví en Madrid.} — presents living there as a completed stretch.
  \item \textbf{Vivía en Madrid.} — puts us inside that past stretch or treats it as background.
\end{itemize}

Both are past. The speaker chooses how to frame the event.
\end{grammarlens}

\begin{cousinweb}
\textbf{Vivir} continues Latin \textbf{vīvere}. Its infinitive belongs to the Spanish \textbf{-ir} class, but the imperfect has converged with \textbf{-er}. That shared row is a reliable fact about these learned forms, not a claim that the two conjugations are identical everywhere.
\end{cousinweb}

\subsection*{Your turn}
Say these aloud:

\begin{itemize}
\raggedright
  \item “vivía, vivías, vivía”
  \item “comía and vivía share -ía”
  \item “Viví en Madrid. Vivía en Madrid.”
\end{itemize}

\begin{tcolorbox}[breakable,colback=teal!4,colframe=teal!35!black,title={Before you move on}]
Give the three \textbf{vivir} forms. (\emph{Vivía, vivías, vivía.}) Which class shares them? (\textbf{-er}.) What contrast can \textbf{viví / vivía} make? (Completed stretch versus an open or background past.) Next: learn \textbf{ver} before asking it to enter this tense.
\end{tcolorbox}
\section[ver]{ver — learn the verb before its past}

\label{lesson:ES-C16-ver}



\begin{quote}
\textbf{Ver} means “to see.” Its written \textbf{v} uses the Spanish sound you already know: it is not contrasted with a separate English-style \textbf{v} sound. Say \textbf{ver}.
\end{quote}

\begin{grammarlens}[title={Grammar Lens: three present forms}]
\noindent
\begin{tabularx}{\linewidth}{@{}>{\raggedright\arraybackslash}X>{\raggedright\arraybackslash}X@{}}
\toprule
\textbf{person} & \textbf{present} \\
\midrule
yo & \textbf{veo} \\
tú & \textbf{ves} \\
él / ella / usted & \textbf{ve} \\
\bottomrule
\end{tabularx}

The \textbf{tú} and third-person forms resemble the familiar \textbf{-er} row. The \textbf{yo} form keeps an extra vowel: \textbf{veo}. Reuse known \textbf{café}: \textbf{Veo café} — “I see coffee.”
\end{grammarlens}

\begin{cousinweb}
Spanish \textbf{ver} continues Latin \textbf{vidēre}, “to see,” through Old Spanish \textbf{veer}. Loss of the consonant between vowels helps explain the shorter shape; the two adjacent vowels later contracted in the infinitive. English \textbf{video} belongs to the same Latin family, a memory bridge rather than a new Spanish word for this lesson.
\end{cousinweb}

\subsection*{Your turn}
Say these aloud:

\begin{itemize}
\raggedright
  \item “ver — to see”
  \item “veo, ves, ve”
  \item “ver comes through veer from vidēre”
\end{itemize}

\begin{tcolorbox}[breakable,colback=teal!4,colframe=teal!35!black,title={Before you move on}]
What does \textbf{ver} mean? (“To see.”) Give its three learned present forms. (\emph{Veo, ves, ve.}) Which Latin verb stands behind it? (\textbf{Vidēre}.) Next: add \textbf{ser} to the open past.
\end{tcolorbox}
\section[era]{era — an old root in an open past}

\label{lesson:ES-C16-ser-imperfecto}



\begin{quote}
You know present \textbf{soy, eres, es} and completed-past \textbf{fui, fuiste, fue}. The open or background past of \textbf{ser} uses a different short stem: \textbf{era}.
\end{quote}

\begin{grammarlens}[title={Grammar Lens: the singular ser row}]
\noindent
\begin{tabularx}{\linewidth}{@{}>{\raggedright\arraybackslash}X>{\raggedright\arraybackslash}X@{}}
\toprule
\textbf{person} & \textbf{open or background past} \\
\midrule
yo & \textbf{era} \\
tú & \textbf{eras} \\
él / ella / usted & \textbf{era} \\
\bottomrule
\end{tabularx}

Say \textbf{era · eras · era}. The first and third forms match, just as they did in the regular imperfect rows. \textbf{¿Era usted?} asks “Was it you?” without adding a new noun. Plural persons wait.
\end{grammarlens}

\begin{cousinweb}
\textbf{Era} continues Latin \textbf{eram}, an old form in the \textbf{esse} family. That is the family already connected to Spanish \textbf{ser} and English \textbf{essence}. The preterite \textbf{fue} comes from another historical set, so \textbf{era / fue} is a real contrast inside one modern verb.
\end{cousinweb}

\subsection*{Your turn}
Say these aloud:

\begin{itemize}
\raggedright
  \item “era, eras, era”
  \item “¿Era usted?”
  \item “era comes from Latin eram”
\end{itemize}

\begin{tcolorbox}[breakable,colback=teal!4,colframe=teal!35!black,title={Before you move on}]
Give the three learned \textbf{ser} forms. (\emph{Era, eras, era.}) Which past form supplies background, \textbf{era} or \textbf{fue}? (\textbf{Era}.) Which Latin form stands behind it? (\textbf{Eram}.) Next: follow \textbf{ir} back toward its own older verb.
\end{tcolorbox}
\section[iba]{iba — ir returns to its inherited past}

\label{lesson:ES-C16-ir-imperfecto}



\begin{quote}
Present \textbf{voy} and completed-past \textbf{fui} came from different historical sets. For an open or habitual past, \textbf{ir} gives you \textbf{iba} — “I was going” or “I used to go.”
\end{quote}

\begin{grammarlens}[title={Grammar Lens: the singular ir row}]
\noindent
\begin{tabularx}{\linewidth}{@{}>{\raggedright\arraybackslash}X>{\raggedright\arraybackslash}X@{}}
\toprule
\textbf{person} & \textbf{open or habitual past} \\
\midrule
yo & \textbf{iba} \\
tú & \textbf{ibas} \\
él / ella / usted & \textbf{iba} \\
\bottomrule
\end{tabularx}

Say \textbf{iba · ibas · iba}. Reuse the destination frame: \textbf{Iba a Madrid} — “I was going to Madrid” or “I used to go to Madrid.” Plural persons wait.
\end{grammarlens}

\begin{cousinweb}
\textbf{Iba} continues Latin \textbf{ībam}, from \textbf{īre}, the inherited verb behind the Spanish infinitive \textbf{ir}. Present \textbf{voy} contains material associated with \textbf{vādere}, while preterite \textbf{fui} comes from the Latin “be” perfect family. These sources help explain the three learned sets; they do not mean that every form in every tense can be assigned to one source by a three-box rule.
\end{cousinweb}

\subsection*{Your turn}
Say these aloud:

\begin{itemize}
\raggedright
  \item “iba, ibas, iba”
  \item “Fui a Madrid. Iba a Madrid.”
  \item “iba continues ībam from īre”
\end{itemize}

\begin{tcolorbox}[breakable,colback=teal!4,colframe=teal!35!black,title={Before you move on}]
Give the three learned forms. (\emph{Iba, ibas, iba.}) Say “I was going to Madrid.” (\emph{Iba a Madrid.}) Which Latin form stands behind \textbf{iba}? (\textbf{Ībam}, from \textbf{īre}.) Next: put newly learned \textbf{ver} into the same past.
\end{tcolorbox}
\section[veía]{veía — keep the extra e you already met}

\label{lesson:ES-C16-ver-imperfecto}



\begin{quote}
You learned \textbf{ver} before its past: \textbf{veo, ves, ve}. Its imperfect keeps \textbf{ve-} and adds the familiar \textbf{-ía} row. That gives \textbf{veía} — “I was seeing” or “I used to see.”
\end{quote}

\begin{grammarlens}[title={Grammar Lens: the singular ver row}]
\noindent
\begin{tabularx}{\linewidth}{@{}>{\raggedright\arraybackslash}X>{\raggedright\arraybackslash}X@{}}
\toprule
\textbf{person} & \textbf{open or habitual past} \\
\midrule
yo & \textbf{veía} \\
tú & \textbf{veías} \\
él / ella / usted & \textbf{veía} \\
\bottomrule
\end{tabularx}

Say \textbf{veía · veías · veía}. Keep \textbf{í-a} in separate syllables. Plural persons wait.
\end{grammarlens}

\begin{cousinweb}
Old Spanish \textbf{veer} displays the \textbf{ve-} shape clearly. The imperfect \textbf{veía} preserves that vowel before the familiar \textbf{-ía} material. It is useful to call this an exceptional learned row, but not to pretend that \textbf{ver} was already known before this chapter or that the history is merely “regular plus one letter.”
\end{cousinweb}

\subsection*{Your turn}
Say these aloud:

\begin{itemize}
\raggedright
  \item “veía, veías, veía”
  \item “Veo café. Veía café.”
  \item “veer helps explain ve-”
\end{itemize}

\begin{tcolorbox}[breakable,colback=teal!4,colframe=teal!35!black,title={Before you move on}]
Give the three learned forms. (\emph{Veía, veías, veía.}) Which vowel stays before \textbf{-ía}? (\textbf{e}.) Which earlier form displays it? (\textbf{Veer}.) Next: retrieve every Chapter 16 row and make a few bounded contrasts.
\end{tcolorbox}
\section[Practice]{Practice — keep the past open or close it}

\label{lesson:ES-C16-practice}



\begin{quote}
This checkpoint adds no vocabulary and no plural persons. Retrieve the seven learned singular rows, then choose whether a known sentence presents the past as completed or leaves it open as background, progress, or habit.
\end{quote}

\begin{grammarlens}[title={Grammar Lens: the complete Chapter 16 frame}]
Regular rows:

\begin{itemize}
\raggedright
  \item \textbf{yo:} \textbf{hablaba · comía · vivía}.
  \item \textbf{tú:} \textbf{hablabas · comías · vivías}.
  \item \textbf{él / ella / usted:} \textbf{hablaba · comía · vivía}.
\end{itemize}

The three special rows:

\begin{itemize}
\raggedright
  \item \textbf{yo:} \textbf{era · iba · veía}.
  \item \textbf{tú:} \textbf{eras · ibas · veías}.
  \item \textbf{él / ella / usted:} \textbf{era · iba · veía}.
\end{itemize}

The first and third forms match in every learned row. \textbf{Comer} and \textbf{vivir} share \textbf{-ía, -ías, -ía}. \textbf{Ver} keeps \textbf{ve-} before that material. Plural persons remain outside this checkpoint.
\end{grammarlens}

\begin{grammarlens}[title={Grammar Lens: one scene, two frames}]
\begin{itemize}
\raggedright
  \item \textbf{Hablé español.} — completed event. \textbf{Hablaba español.} — open, background, or habitual speaking.
  \item \textbf{Viví en Madrid.} — completed stretch. \textbf{Vivía en Madrid.} — we stand inside that past stretch.
  \item \textbf{Fui a Madrid.} — completed trip. \textbf{Iba a Madrid.} — motion was in progress or habitual.
\end{itemize}

The tense does not mechanically report how long something lasted. It frames the past from a different viewpoint.
\end{grammarlens}

\begin{cousinweb}
\textbf{-aba} continues the Latin \textbf{-ābam} family. \textbf{-ía} continues material from the \textbf{-ēbam} family after several Romance changes. \textbf{Era} continues \textbf{eram}; \textbf{iba} continues \textbf{ībam} from \textbf{īre}; and \textbf{veía} keeps the \textbf{ve-} shape also visible in Old Spanish \textbf{veer}. These are historical bridges, not recipes for generating unlearned forms.
\end{cousinweb}

\subsection*{Your turn}
Say these aloud:

\begin{itemize}
\raggedright
  \item “hablaba, comía, vivía”
  \item “era, iba, veía”
  \item “Hablé / hablaba · viví / vivía · fui / iba”
\end{itemize}

\begin{tcolorbox}[breakable,colback=teal!4,colframe=teal!35!black,title={Before you move on}]
Give the six \textbf{yo} forms. (\emph{Hablaba, comía, vivía, era, iba, veía.}) Which two regular classes share a row? (\textbf{-er} and \textbf{-ir}.) What is the core contrast? (The preterite presents a completed event; the imperfect leaves past action open, habitual, in progress, or in the background.) Chapter 17 can now build on \textbf{-ía} without borrowing a plural form here.
\end{tcolorbox}
